\section{Introduction}
\label{sec:introduction}

Agentic AI systems are rapidly evolving from passive question-answering assistants into autonomous decision-makers that plan, call tools, and act in open-ended environments~\cite{mialon2023gaia,drouin2024workarena,xie2024osworld}. This shift is especially visible in information-retrieval (IR) contexts, where modern agents now act on retrieved documents and structured data---issuing follow-up queries, writing files, sending messages, and modifying records---rather than merely returning passages~\cite{xi2024agentir}. As these systems become embedded in real-world workflows, their failures are no longer limited to incorrect answers on static tasks. Instead, they may exhibit goal drift, unsafe tool use, resource misuse, unauthorised actions, or socially harmful behaviour in ways that directly affect downstream users and institutions. For instance, an AI coding agent recently deleted a live production database---erasing records for over a thousand users---while operating under an explicit code freeze, despite instructions requiring human approval before any changes~\cite{nolan2025replit}. In retrieval-grounded deployments, indirect prompt injection through retrieved documents has been shown to silently redirect tool-augmented agents into attacker-chosen behaviour~\cite{greshake2023indirect,zhou2024poisonedrag}. Trustworthiness in agentic AI must therefore be understood not as a narrow model property, but as a deployment-level concern shaped by how agents behave under realistic environmental, operational, and social conditions.

Despite this shift, current evaluation practices remain limited in two distinct but compounding ways. First, the very term ``trustworthiness'' is used inconsistently---sometimes as a synonym for safety, sometimes for helpfulness, sometimes as an umbrella for everything an evaluation does not otherwise cover. Existing surveys list many desiderata~\cite{liu2024trustllm,wang2024trustworthyagent} but do not give a small, measurable set of properties against which an agentic system can be assessed; without such a definition, any claim that an agent ``passes a trustworthiness audit'' is itself vacuous. Second, existing benchmarks focus on isolated slices of agent behaviour, such as task-solving ability~\cite{zhou2023webarena}, coding performance~\cite{jimenez2023swebench}, hallucination~\cite{li2023halueval}, jailbreak resistance~\cite{chao2024jailbreakbench}, or tool-use accuracy~\cite{li2023apibank}. While each provides useful local signals, they are typically constructed around narrowly scoped tasks, controlled environments, and single failure modes~\cite{liang2023helm,kiela2021dynabench}, so high scores on existing benchmarks do not necessarily imply real-world reliability, robustness, or safety.

We argue that the central limitation is therefore not merely that current evaluations cover too few dimensions, but that they lack \emph{both} (i)~a concrete definition of trustworthiness and (ii)~a principled notion of \emph{representativeness}. Agent trustworthiness is not a point property that can be inferred from a handful of benchmark instances; rather, it is a distributional property that depends on how an agent behaves over a heterogeneous, heavy-tailed space of socio-technical scenarios. When the distribution of benchmark tasks is misaligned with the distribution of real deployment conditions, evaluation results can create a false sense of confidence~\cite{koh2021wilds,shirali2022dynamicbench}. This problem is particularly severe for low-frequency but high-consequence events, including adversarial manipulations, cascading tool failures, and socially sensitive interactions, which are often absent from standard evaluation suites despite being central to trustworthy deployment.

To address both gaps, we propose the \textbf{Holographic Agent Assessment Framework (HAAF)}---a systematic trustworthiness-evaluation framework for agentic AI.
Our key idea is to first \emph{define} agent trustworthiness as a five-property profile (defined in \S\ref{subsec:defining_trustworthiness}) and then to \emph{measure} that profile over a \emph{scenario manifold} that spans task types, tool interfaces, interaction dynamics, social contexts, and risk levels, rather than over disconnected benchmark islands. The framework centres on a \emph{distribution-aware representative sampling engine} that defines \emph{what} should be tested, while three complementary probing interfaces---static cognitive and policy analysis, interactive sandbox simulation, and social-ethical alignment assessment---together define \emph{how} agents \emph{should} be evaluated; the present instantiation (\S\ref{sec:experiments}) exercises the interactive sandbox layer in depth, with the static and social-ethical layers described as part of the framework but reserved for future cycles. These components are connected through an iterative \emph{Trustworthy Optimization Factory}: red-team probing exposes vulnerability surfaces, blue-team hardening designs targeted interventions, and re-evaluation verifies improvement, cycling until the deployed system meets deployment-readiness thresholds. Rather than evaluating model weights in isolation, we assess \emph{deployed agentic systems} comprising the model, prompting policies, tool interfaces, guardrails, and oversight mechanisms.

We illustrate this paradigm with concrete evidence: a coverage audit of existing benchmarks across eight deployment-relevant evaluation dimensions that map onto the five properties (Table~\ref{tab:coverage}); a complete red-team/blue-team cycle on a single focal agent that reduces its violation rate from 16.7\% to 4.2\% with no new improper refusals; a \emph{cross-model trustworthiness profile} over 13 contemporary agentic systems spanning seven model families (Llama, Mistral, Kimi, GLM, Qwen, GPT-oss, DeepSeek) served through both local and Amazon Bedrock back-ends; and a \emph{cross-family generalisation experiment} in which the focal-model interventions are applied uniformly to those same 13 systems and observed to transfer in 12 of 13 cases (one intervention-resistant outlier). The cross-model profile surfaces property-level trade-offs---one system stronger on robustness, another on social-ethical alignment---that any single aggregate metric would erase.

\paragraph{Representativeness: target vs.\ current evidence.}
We separate the \emph{target} of representativeness from the \emph{evidence} we currently offer for it. Our 100 scenarios (s01--s100) are hand-authored: 24 design-study scenarios (s01--s24) drove the focal-model Factory cycle, and 76 validation scenarios (s25--s100) instantiate the five-property profile over the remaining HAAF cells. The five sandbox tools (\texttt{search\_docs}, \texttt{db\_query}, \texttt{read\_file}, \texttt{write\_file}, \texttt{send\_message}) were chosen so that each maps to a top-$K$ operation class observed in public-trace agent benchmarks (SWE-bench~\cite{jimenez2023swebench}, API-Bank~\cite{li2023apibank}, WebArena~\cite{zhou2023webarena}, AgentBench~\cite{liu2023agentbench}). We therefore claim the \emph{target}---a representative distribution over deployment-relevant operation classes---and the \emph{methodology} (Layer-4 sampling objective with a deployment-grounded weighting hook), but \emph{not} that the present 100 scenarios constitute a statistically representative slice of production traffic. Grounding scenario weights and tool mixes in real telemetry is committed as v2 (\S\ref{sec:conclusion}).

\paragraph{Contributions.}
In summary, this paper makes the following six contributions:
\begin{enumerate}[leftmargin=*,label=(C\arabic*)]
\item \textbf{Definition.} We give an operational five-property definition of agentic trustworthiness (P1--P5), grounded in current AI risk frameworks~\cite{nist2023airmf,eu2024aiact}, in which \emph{Improper Refusal} is a first-class taxonomy class mapped to P1 rather than a side metric.
\item \textbf{Cross-model evaluation.} We evaluate 13 contemporary models across seven families (Llama, Mistral, Kimi, GLM, Qwen, GPT-oss, DeepSeek) on a 100-scenario validation suite, producing $2{,}600$ Control/Treated trajectories---substantially broader than the 24-scenario, single-focal-model design study that drove the Factory cycle.
\item \textbf{End-to-end instantiation.} We instantiate the framework on a single shared sandbox: red-team attribution of property-level failures, blue-team interventions targeted at those failures, and re-evaluation showing the interventions reduce $\mathrm{RWF}$ on all 13 systems without per-model or per-scenario tuning---12 substantially, while one system (Qwen3-32B) is intervention-resistant (marginal $\Delta{=}+0.019$).
\item \textbf{Layer-4 reference implementation.} We release a reference implementation of the Layer-4 distribution-aware sampling objective; a 40-scenario subset it selects preserves the cross-model ranking against the full 100-scenario suite (Kendall-$\tau$\,=\,0.890, Spearman-$\rho$\,=\,0.963).
\item \textbf{Statistical confidence.} We report Wilson 95\%~CIs on per-property success rates and paired-scenario bootstrap ($B{=}10{,}000$) on four intra-family anti-scaling pairs; two of four clearly reject the larger-is-no-worse null at $p<0.05$, with a third sitting at the boundary ($p=0.050$) (with $n{=}14$--$26$ scenarios per property, per-cell CIs span roughly $\pm 0.2$, so these statistics are descriptive of cross-family trends rather than powered per-property claims; see \S\ref{sec:discussion}).
\item \textbf{Explicit scope.} We are explicit about what the current artefact is not: scenarios and sandbox tools are hand-authored, IR validation uses a single LLM-judge (Kimi-K2.5, itself one of the 13 evaluated agents---a deliberate scope limitation discussed in \S\ref{sec:discussion}), and scenario weights are not yet grounded in production traces. A deployment-grounded sampler is committed as v2.
\end{enumerate}
